\documentclass[11pt, letterpaper]{article}

% --- Packages ---
\usepackage{amsmath}
\usepackage{graphicx}
\usepackage{amsfonts}
\usepackage{enumitem}
\usepackage{xcolor}
\usepackage{listings}
\usepackage{tcolorbox}
\usepackage[left=1in, right=1in, bottom=1in, top=3cm, headheight=2cm, headsep=0.5cm, footskip=1cm]{geometry}
\usepackage{fancyhdr}

% --- Code Listing Formatting ---
\definecolor{codegreen}{rgb}{0,0.6,0}
\definecolor{codegray}{rgb}{0.5,0.5,0.5}
\definecolor{codepurple}{rgb}{0.58,0,0.82}
\definecolor{backcolour}{rgb}{0.95,0.95,0.92}

\lstset{ 
    backgroundcolor=\color{backcolour},   
    commentstyle=\color{codegreen},
    keywordstyle=\color{blue},
    numberstyle=\tiny\color{codegray},
    stringstyle=\color{codepurple},
    basicstyle=\ttfamily\footnotesize,
    breakatwhitespace=false,         
    breaklines=true,                 
    captionpos=b,                    
    keepspaces=true,                 
    numbers=left,                    
    numbersep=5pt,                  
    showspaces=false,                
    showstringspaces=false,
    showtabs=false,                  
    tabsize=2,
    language=C++
}

% --- Header Configuration ---
\pagestyle{fancy} 
\fancyhf{} 
\renewcommand{\headrulewidth}{0.4pt} 
\lhead{\raisebox{-0.5cm}{\textbf{Universidad de Antioquia}}}
\rhead{\large \textbf{Code Walkthrough: Mass Optimization} \\ \normalsize Spaceflight Dynamics 2026-1} 
\cfoot{\thepage} 

\begin{document}

\begin{center}
    {\Large \textbf{Deconstructing the Interplanetary Optimization Script}} \\
    \vspace{0.2cm}
    {\large \textit{Connecting SQP Mathematics to GMAT Syntax}} \\
    \vspace{0.4cm}
    \textit{Prof. Juan | Aerospace Engineering Program}
\end{center}
\hrule
\vspace{0.5cm}

\section*{Script Anatomy: Multi-Burn Mass Optimization}
In the previous lab, we introduced a script template designed to optimize an interplanetary transfer using Fuel Mass rather than squared $\Delta V$. This walkthrough explains the underlying astrodynamics and mathematical principles driving each block of that script.

% ---------------------------------------------------------
\subsection*{1. Hardware Association (The Physics Engine)}
\begin{lstlisting}[numbers=none]
Tank1.FuelMass = 500.0;
DepartureBurn.Tank = {Tank1};
ArrivalBurn.Tank = {Tank1};
\end{lstlisting}
\textbf{The Concept:} When we used basic Differential Correctors (Targeters), we did not care about fuel tanks. A targeter simply calculates the kinematic $\Delta V$ vector required to intersect a point in space. 
\\
\textbf{Why it changes here:} Because our optimization strategy relies on the Rocket Equation (exploiting the Oberth Effect), GMAT must be able to physically deplete mass from the spacecraft during the \texttt{Maneuver} command. By linking both burns to \texttt{Tank1}, GMAT continuously integrates the mass drop, dynamically changing the spacecraft's acceleration and inertia.

% ---------------------------------------------------------
\subsection*{2. The Control Variables (The ``Vary'' Commands)}
\begin{lstlisting}[numbers=none]
Vary Yukon1(DepartureBurn.Element1 = 2.0, {Perturbation = 0.0001, MaxStep = 0.5});
Vary Yukon1(ArrivalBurn.Element1 = 1.0, {Perturbation = 0.0001, MaxStep = 0.5});
\end{lstlisting}
\textbf{The Concept:} We are giving the Yukon (SQP) algorithm two ``knobs'' to turn: the prograde velocity of the Earth departure and the prograde velocity of the Mars arrival. 
\\
\textbf{The Parameters:}
\begin{itemize}[noitemsep]
    \item \textbf{Initial Guess (2.0 and 1.0):} Gradient descent algorithms are blind. If you guess 0.0, the solver might fall into a mathematical flatland and fail. You must provide a reasonable physical estimate.
    \item \textbf{Perturbation (0.0001):} This is the ``finite differencing'' step. Yukon will invisibly add $0.0001$ km/s to your burn, run the whole simulation, and see how the final radius changes to calculate the gradient (Jacobian matrix).
    \item \textbf{MaxStep (0.5):} This prevents the algorithm from making wild, unstable leaps across the mathematical landscape during early iterations.
\end{itemize}

% ---------------------------------------------------------
\subsection*{3. The Sequence Flow (Propagate vs. Maneuver)}
\begin{lstlisting}[numbers=none]
Propagate Prop1(Sat) {Sat.Earth.Periapsis};
Maneuver DepartureBurn(Sat);
Propagate Prop1(Sat) {Sat.Earth.Apoapsis};
Maneuver ArrivalBurn(Sat);
\end{lstlisting}
\textbf{The Concept:} Order is absolute in a Solver Control Sequence. The optimizer executes this chronologically. We purposely force the \texttt{DepartureBurn} to occur exactly at \texttt{Periapsis} to maximize the spacecraft's velocity when the engine fires, deliberately forcing the optimizer to utilize the \textbf{Oberth Effect}. 

% ---------------------------------------------------------
\subsection*{4. The Constraint and the Trap (The Objective Function)}

\begin{tcolorbox}[colback=red!5!white,colframe=red!75!black,title=The Optimization Trap: What are you actually minimizing?]
In mathematics, if you tell a computer to \texttt{Minimize(FuelMass)}, the computer will gleefully burn your engines until the tank is empty, driving the remaining mass to zero! We want to minimize the fuel \textbf{spent}, which means maximizing the fuel \textbf{remaining}.
\end{tcolorbox}

\begin{lstlisting}[numbers=none]
% 1. The Fence
NonLinearConstraint Yukon1(Sat.Earth.Apoapsis = 150000.0);

% 2. The Objective
Create Variable FuelSpent;
FuelSpent = 500.0 - Sat.FuelTank1.FuelMass;
Minimize Yukon1(FuelSpent);
\end{lstlisting}

\textbf{The Concept:}
\begin{itemize}
    \item \textbf{The Constraint:} The \texttt{NonLinearConstraint} acts as the Active-Set boundary. The solver will adjust the \texttt{Vary} commands until the apoapsis hits exactly 150,000 km.
    \item \textbf{The Objective:} To avoid the trap mentioned above, we create a custom variable called \texttt{FuelSpent}. By defining it as $(\text{Initial Mass} - \text{Current Mass})$, we give the Yukon solver a smooth, strictly positive mathematical bowl to minimize. 
\end{itemize}

\textbf{Summary Conclusion:} 
Because the fuel spent scales exponentially with $\Delta V$ (via the Rocket Equation), Yukon realizes that adding 1 m/s at the Arrival Burn costs more mass than adding 1 m/s at the Departure Burn (due to the Oberth effect). Therefore, the SQP algorithm will naturally pile as much $\Delta V$ into \texttt{DepartureBurn} as physically possible without splitting the burns evenly!

\end{document}